\section[Теоретические сведения об эллиптических кривых]{Теоретические сведения об эллиптических кривых\ над простыми конечными полями}

В работе рассматриваются эллиптические кривые над простыми конечными полями $\mathbb{F}_p$, где $p$ --- нечётное простое число. Такой выбор соответствует постановке исследуемой задачи и охватывает класс кривых, на котором проводились вычислительные эксперименты. При этом внимание сосредоточено не на программной реализации арифметики точек, а на тех математических свойствах кривых, которые необходимы для формулировки задачи дискретного логарифмирования и анализа алгоритмов её решения.

Под конечным простым полем $\mathbb{F}_p$ понимается поле вычетов по простому модулю $p$. Его элементы можно отождествить с числами $0,1,\dots,p-1$, а операции сложения и умножения выполняются по модулю $p$. Простота модуля гарантирует, что каждый ненулевой элемент имеет мультипликативно обратный, поэтому $\mathbb{F}_p$ действительно является полем. Именно над такими полями удобно рассматривать эллиптические кривые в короткой форме Вейерштрасса.

Эллиптическая кривая над полем $\mathbb{F}_p$ задаётся уравнением
\[
E: y^2 = x^3 + ax + b,
\]
где $a,b \in \mathbb{F}_p$. Однако не всякое такое уравнение определяет допустимую эллиптическую кривую. Необходимо выполнение условия невырожденности
\[
4a^3 + 27b^2 \not\equiv 0 \pmod p.
\]
Если это условие нарушается, кривая имеет особые точки, и на множестве её решений уже нельзя корректно определить ту групповую структуру, которая используется в криптографии. Поэтому невырожденность является базовым требованием при выборе параметров.

Множество точек кривой над полем $\mathbb{F}_p$ состоит из всех пар $(x,y) \in \mathbb{F}_p \times \mathbb{F}_p$, удовлетворяющих уравнению кривой, а также специальной точки на бесконечности $\mathcal{O}$. Это множество обозначается через $E(\mathbb{F}_p)$. Включение точки $\mathcal{O}$ принципиально важно: именно она играет роль нейтрального элемента группы. Таким образом,
\[
E(\mathbb{F}_p)=\{(x,y)\in \mathbb{F}_p^2 \mid y^2=x^3+ax+b\}\cup\{\mathcal{O}\}.
\]
Так как поле конечно, множество $E(\mathbb{F}_p)$ тоже конечно.

На множестве $E(\mathbb{F}_p)$ определяется абелева групповая операция сложения точек. Геометрически над полем вещественных чисел эта операция интерпретируется через проведение прямой через две точки кривой и отражение полученной третьей точки пересечения относительно оси абсцисс. Над конечным полем сохраняются те же алгебраические формулы, хотя сама геометрическая картина уже не используется буквально. Для дальнейшего изложения существенно не детальное устройство формул сложения и удвоения, а тот факт, что множество точек образует конечную абелеву группу.

Наличие групповой структуры позволяет ввести важнейшую операцию --- умножение точки на целое число. Для точки $P\in E(\mathbb{F}_p)$ и целого $k\ge 0$ величина $kP$ определяется как сумма $P$ с самой собой $k$ раз, причём по определению $0P=\mathcal{O}$. Скалярное умножение является основной операцией в группе точек эллиптической кривой.

Если существует наименьшее положительное число $n$, для которого выполнено равенство
\[
nP=\mathcal{O},
\]
то число $n$ называется порядком точки $P$. Тогда множество
\[
\langle P\rangle=\{\mathcal{O},P,2P,\dots,(n-1)P\}
\]
образует циклическую подгруппу группы $E(\mathbb{F}_p)$ порядка $n$, порождённую точкой $P$.

Помимо порядка отдельной точки рассматривается и порядок всей кривой, то есть число точек множества $E(\mathbb{F}_p)$. Оно обозначается $\#E(\mathbb{F}_p)$. Для этого числа справедлива оценка Хассе
\[
\left|\#E(\mathbb{F}_p)-(p+1)\right|\le 2\sqrt{p}.
\]
Следовательно, число точек на кривой всегда близко к $p+1$, однако точное значение зависит от коэффициентов $a$ и $b$.

С точки зрения структуры конечная абелева группа точек эллиптической кривой над $\mathbb{F}_p$ изоморфна произведению не более чем двух циклических групп. В частности, не вся группа $E(\mathbb{F}_p)$ обязана быть циклической. Поэтому в дальнейшем важно различать всю группу точек кривой и конкретные циклические подгруппы, порождённые отдельными точками.